\documentclass[12pt,a4paper]{article}

\usepackage[margin=2.5cm]{geometry}
\usepackage{graphicx}
\usepackage{booktabs}
\usepackage{array}
\usepackage{amsmath}
\usepackage{siunitx}
\usepackage{natbib}
\usepackage{caption}
\usepackage{subcaption}
\usepackage{placeins}
\usepackage{setspace}
\usepackage{lineno}
\usepackage{url}
\usepackage[hidelinks]{hyperref}

\graphicspath{
  {figures/}
  {D:/Data_Analysis_R/Gas_Concentration_Emission_R/workflows/mapping_high_resolution/plots/manuscript1/}
}

\sisetup{
  separate-uncertainty = true,
  per-mode = symbol,
  range-phrase = {--},
  range-units = single
}

\onehalfspacing
\linenumbers

\title{Spatial heterogeneity of carbon dioxide, methane, and ammonia in a naturally ventilated dairy barn: high-resolution mapping and implications for sampling design}

\author{
Harsh Sahu$^{a,b}$,
Long Chen$^{a}$,
Thomas Amon$^{a,b}$,
David Janke$^{a,*}$
}

\date{}

\begin{document}

\maketitle

\begin{center}
\small
$^{a}$Leibniz Institute for Agricultural Engineering and Bioeconomy (ATB),
Max-Eyth-Allee 100, 14469 Potsdam, Germany\\
$^{b}$Institute of Animal Hygiene and Environmental Health, Department of
Veterinary Medicine, Freie Universit\"at Berlin, Germany\\
$^{*}$Corresponding author: \href{mailto:djanke@atb-potsdam.de}{djanke@atb-potsdam.de}
\end{center}

\begin{abstract}
Homogeneous gas mixing is commonly assumed when concentrations are used to
characterise air quality or calculate emissions from naturally ventilated
dairy barns, but the magnitude and spatial structure of departures from this
assumption remain poorly resolved. Carbon dioxide (\(\mathrm{CO_2}\)),
methane (\(\mathrm{CH_4}\)), and ammonia (\(\mathrm{NH_3}\)) concentrations
were therefore mapped at 51 internal locations arranged as 17 vertical
triplets. Two cavity ring-down spectrometers and two Fourier-transform infrared
spectrometers were operated at a temporal resolution of four minutes; the
Fourier-transform infrared measurements were harmonised to the cavity
ring-down spectroscopy scale using a preliminary 24 h co-located comparison.
All gases exhibited structured vertical and horizontal heterogeneity. Mean
concentrations decreased from top to middle to bottom, but the magnitude of
the vertical gradient differed among gases and horizontal positions. The
\(\mathrm{CH_4}/\mathrm{CO_2}\) and
\(\mathrm{NH_3}/\mathrm{CO_2}\) ratios also decreased with height group from
top to bottom, showing that the observed structure was not explained solely by
uniform dilution or concentration. As a practical consequence of this spatial
analysis, alternative height configurations were evaluated in 743 complete
two-hour blocks. Top-only and bottom-only measurements produced systematic,
gas-specific biases, whereas a spatially distributed top-and-bottom network
reproduced the full-density mean with Lin's concordance correlation
coefficients of 0.999 and mean relative biases below 0.4\%. The primary finding
is that gas concentrations within the barn were spatially structured rather
than homogeneous; understanding this structure subsequently enabled the
network to be reduced from 51 to 34 locations without materially changing the
barn mean.
\end{abstract}

\noindent\textbf{Keywords:}
air quality; dairy housing; horizontal gradient; mixing ratio; spatial
variability; vertical profile

\section*{Nomenclature}

\begin{tabular}{>{\raggedright\arraybackslash}p{2.2cm}
                >{\raggedright\arraybackslash}p{10.8cm}}
\toprule
Symbol & Definition\\
\midrule
\(\mathrm{CH_4}\) & methane\\
\(\mathrm{CO_2}\) & carbon dioxide\\
CCC & Lin's concordance correlation coefficient\\
CRDS & cavity ring-down spectroscopy\\
FTIR & Fourier-transform infrared spectroscopy\\
\(\mathrm{H_2O}\) & water vapour\\
\(\mathrm{NH_3}\) & ammonia\\
NVDB & naturally ventilated dairy barn\\
RE & relative deviation from the full-density estimate, \%\\
RMSE & root mean square error\\
\(s\) & south-background sampling location\\
SE & standard error\\
SP & sampling point\\
\bottomrule
\end{tabular}

\section{Introduction}

Emission estimates for naturally ventilated livestock buildings depend on
both the air exchange rate and the concentration difference across the
measurement boundary. Neither component is readily constrained because wind
and buoyancy can reverse flow through large openings, while animals, manure
surfaces, internal obstructions, and mechanical mixing devices generate local
sources and transport pathways. Uncertainty in a concentration measurement is
therefore not only analytical; it also includes whether the selected location
represents the air volume or control surface of interest. These distinctions
are important because a representative internal concentration is not
necessarily representative of an emission flux when the corresponding airflow
field is unknown \citep{Calvet2013,Ogink2013}.

The consequences of imperfect mixing have been demonstrated under controlled
and field conditions. In a mechanically ventilated test installation,
ventilation-rate errors of up to 86\% arose when tracer concentrations from
individual internal positions were used, although this value describes a
tracer-gas ventilation estimate rather than a barn concentration or emission
error \citep{VanBuggenhout2009}. Wind-tunnel measurements of a large naturally
ventilated building likewise showed that emission estimates based on
concentration and velocity at one constant height could differ markedly from
vertically resolved measurements \citep{Janke2020}. These studies establish
that the spatial arrangement of measurements can be as important as their
nominal precision, but their controlled flow fields do not reproduce the full
temporal variability of an occupied dairy barn.

Field studies have identified heterogeneity in both horizontal and vertical
directions. Height-dependent differences in \(\mathrm{NH_3}\),
\(\mathrm{CO_2}\), and tracer-gas mixing ratios were observed in a naturally
ventilated dairy cow barn, with more consistent mixing above the
animal-occupied zone \citep{Mendes2015}. Gas-specific vertical profiles and
incomplete correlations among heights were also reported in an open-sided
dairy barn \citep{Durso2024}. Horizontally, the mean and its variability were
sensitive to the number and separation of sampling locations
\citep{Durso2022}. External wind speed and direction further altered internal
concentrations and air exchange in a site-specific manner
\citep{Saha2013}. Consequently, a vertical profile observed at one horizontal
position cannot be assumed to apply throughout a barn, nor can the behaviour
of one gas be transferred directly to another.

Spatial representativeness must also be tested across contrasting temporal
regimes. In naturally ventilated buildings, the pressure distribution around
the envelope changes with wind speed and approach direction, whereas the
relative contribution of buoyancy changes with the indoor--outdoor temperature
difference. The resulting ventilation regime can alter both the magnitude and
the spatial ordering of sampling-point concentrations \citep{Saha2013}.
Measurements made during one meteorological period therefore cannot establish
whether a spatial pattern is structurally persistent or conditional on the
prevailing weather. Repeated dense-network observations in summer and autumn,
followed by a reduced-network application in late autumn and winter, provide a
means of examining that persistence. Such comparisons require concurrent
meteorological information and careful separation of season from campaign,
analyser deployment, and sampling design; regional observations can describe
the synoptic setting but cannot reproduce the barn-scale airflow field.

Gas ratios provide an additional view of mixing because responses shared by
two gases can partly cancel, whereas differences in source position,
reactivity, or transport remain. Spatial variability in ammonia-to-tracer
mixing ratios has shown that apparent agreement at one height can coexist with
horizontal differences elsewhere in the building \citep{Mendes2015}.
Simulation-based work has consequently investigated representative sampling
heights \citep{Doumbia2021}, and information-based approaches have been used
to identify non-redundant sensor locations across simulated ventilation cases
\citep{Chen2026}. However, the latter approach was developed for
CFD-derived local mean age of air and related flow variables; its direct
transfer to sequential field measurements of gas concentration has not been
established.

Dense field sampling can characterise these patterns but increases tubing,
switching-cycle duration, maintenance requirements, and the risk of sampling
line failure. Conversely, reduction based only on sensor count or a geometric
centre may discard vertical extremes or persistent local effects. Sampling
density and spatial arrangement have both affected airflow estimates in a
naturally ventilated mock-up building \citep{DeVogeleer2017}, yet evidence is
limited on how a full-scale, three-dimensional concentration network can be
reduced while retaining both its dense-network mean and meaningful spatial
structure over extended field campaigns.

Accordingly, three questions were addressed. First, how large and persistent
is the spatial heterogeneity of \(\mathrm{CO_2}\), \(\mathrm{CH_4}\),
\(\mathrm{NH_3}\), and their ratios across height, horizontal position, and
temporal aggregation? It was hypothesised that location and height effects
would be gas-specific and would persist most strongly for \(\mathrm{NH_3}\)
and nitrogen-containing ratios. Second, can the middle locations be omitted
while retaining both the dense-network concentration estimate and the
observed vertical structure? It was hypothesised that a spatially distributed
top-and-bottom design would agree more closely with the 51-location reference
than either extreme alone while preserving vertical contrast. Third, can one
physically eligible location provide a pragmatic approximation of the network
mean? It was hypothesised that a stable candidate could be ranked, but that
gas- and campaign-specific residual bias would preclude universal reference
status.

\section{Materials and methods}

\subsection{Experimental barn}

Measurements were conducted in the experimental dairy barn of the
Lehr- und Versuchsanstalt f\"ur Tierzucht und Tierhaltung, Gross Kreutz,
Brandenburg, Germany. The naturally ventilated barn was \SI{38.88}{m} long and
\SI{17.65}{m} wide. The asymmetrical fibre-cement roof reached approximately
\SI{6.2}{m} at the ridge and approximately \SI{3.6}{m} at the sidewalls. The
internal airspace was approximately \SI{4529}{m^3}. The loose-housing barn was
designed for 58 dairy cows and contained littered cubicles, concrete
walkways, a manure scraper, and adjacent milking and young-stock areas.

\subsection{Sampling layout}

Fifty-one internal SPs were distributed over 17 horizontal positions
(Fig.~\ref{fig:layout}). Three vertically aligned locations were installed at
each position: a top location \SI{60}{cm} below the roof, a middle location
\SI{360}{cm} above the floor, and a bottom location \SI{260}{cm} above the
floor. The location numbers were assigned consecutively within each triplet.
Thus, locations \(1,4,\ldots,49\) represented the top group; locations
\(2,5,\ldots,50\) represented the middle group; and locations
\(3,6,\ldots,51\) represented the bottom group. A single south-background SP,
denoted \(s\), was treated conceptually as location 52 but was excluded from
all analyses of internal spatial heterogeneity.

\begin{figure}[htbp]
\centering
\begin{subfigure}{\textwidth}
  \centering
  \includegraphics[width=\textwidth]{Fig1a_campaign1_sampling_layout_no_ringline.png}
  \caption{Plan view.}
\end{subfigure}

\vspace{0.5em}
\begin{subfigure}{0.92\textwidth}
  \centering
  \includegraphics[width=\textwidth]{Fig1b_barn_cross_section.png}
  \caption{Cross-sectional representation of the sampling heights and fans.}
\end{subfigure}
\caption{Experimental barn and high-resolution sampling design. Orange symbols
indicate locations \SI{60}{cm} below the roof, green symbols indicate
locations \SI{360}{cm} above the floor, and blue symbols indicate locations
\SI{260}{cm} above the floor. The numbered internal locations comprise 17
vertical triplets. The yellow point labelled Background represents the
south-background location \(s\), which was retained in the dataset but
excluded from internal spatial comparisons.}
\label{fig:layout}
\end{figure}

\subsection{Gas analysers and preliminary field comparison}

Concentrations of \(\mathrm{CO_2}\), \(\mathrm{CH_4}\),
\(\mathrm{NH_3}\), and water vapour were measured using two CRDS analysers
(Picarro G2508; Picarro Inc., Santa Clara, CA, USA; CRDS8 and CRDS9) and two
FTIR analysers (Gasmet CX4000; Gasmet Technologies Oy, Vantaa, Finland; FTIR1
and FTIR2). Before the first
mapping campaign, all four analysers were operated for 24 h at a co-located
central barn position. Two polytetrafluoroethylene sampling tubes were used:
one was divided with a T-piece between the two CRDS analysers and the other
was divided between the two FTIR analysers. The FTIR instruments measured
\(\mathrm{CO_2}\) and \(\mathrm{CH_4}\) concentrations that were 6\% higher
than those measured by CRDS, while FTIR \(\mathrm{NH_3}\) was 9\% higher.
Accordingly, FTIR measurements from Campaign 1 were harmonised to the CRDS
scale as

\begin{align}
C_{\mathrm{CO_2,corr}} &= C_{\mathrm{CO_2,raw}}/1.06,\\
C_{\mathrm{CH_4,corr}} &= C_{\mathrm{CH_4,raw}}/1.06,\\
C_{\mathrm{NH_3,corr}} &= C_{\mathrm{NH_3,raw}}/1.09.
\end{align}

CRDS measurements and all water-vapour measurements were retained without
correction. Both raw and corrected concentrations were preserved in the
analytical dataset. Before and after the field comparison, all analysers were
tested with certified cylinders containing approximately \SI{700}{ppm}
\(\mathrm{CO_2}\), \SI{7}{ppm} \(\mathrm{CH_4}\), and \SI{5}{ppm}
\(\mathrm{NH_3}\); deviations remained within the manufacturer-stated
\(\pm 5\%\) accuracy. The field comparison was therefore used for
cross-platform alignment under representative sampling conditions. The CRDS
instruments were treated as the operational comparison scale, not as an
absolute metrological reference. Consequently, the correction removed the
observed mean platform offset but could not separate instrument-specific
drift from spatial effects during the subsequent fixed analyser assignments.

\subsection{Measurement campaigns and data processing}

Campaign 1 data were recovered for June--August 2024. Concurrent ordered
cycling by all four analysers was identified from 3 June to 16 July and again
from 1 to 26 August, with shorter interruptions retained as gaps. CRDS9
sampled locations 1--16, FTIR2 sampled locations 17--26,
CRDS8 sampled locations 27--42, and FTIR1 sampled locations 43--51 and \(s\).
The resulting network comprised 51 internal locations and the
south-background location. Gas measurements were recorded at a four-minute
temporal resolution. Each settled CRDS valve step was identified from the
ordered 1--16 sequence. The transition signal was excluded, the first
\SI{60}{s} was discarded, and the subsequent \SI{180}{s} was averaged.
Across June--August, CRDS9 contained 22,384 valid four-minute steps and two
isolated three-minute steps; CRDS8 contained 26,099 four-minute steps and one
three-minute step. For every three-minute step, \SI{120}{s} remained after
the common \SI{60}{s} flushing period. Single-position and irregular
sequences were retained in audit tables but excluded from the campaign data.
FTIR exports were cleaned by removing repeated embedded headers and retaining
the required timestamp, sampling-position, and gas fields. Their instrument
timestamps were retained without an additional flushing operation. All
timestamps were parsed in the Europe/Berlin time zone, duplicate combinations
of timestamp, analyser, and location were removed within each campaign, and
the retained rows were ordered by timestamp, analyser, and location.

Campaign 2 was conducted from 16 November 2024 at 00:06:47 to 31 December
2024 at 23:57:10 using CRDS8 and CRDS9. The 17 middle locations were omitted
following the Campaign 1 density assessment. Locations 19 and 40 were also
excluded because of fan proximity and a sampling-line leakage problem,
respectively, resulting in 32 internal locations. Each CRDS sampling step
lasted \SI{240}{s}; the first \SI{60}{s} was discarded for flushing and the
remaining \SI{180}{s} was averaged.

Non-positive gas records were treated as missing in the manuscript analytical
dataset. Carbon dioxide values below \SI{300}{ppm} were also treated as
invalid transition or analyser records because they were physically
implausible for barn or ambient air. No row was deleted from the immutable
source dataset.

\begin{table}[htbp]
\centering
\caption{Measurement design and analytical records retained for the two
campaigns. The south-background location \(s\) was retained in the data but
excluded from the count of internal locations.}
\label{tab:campaigns}
\begingroup
\small
\setlength{\tabcolsep}{1.5pt}
\begin{tabular}{@{}c l
  >{\raggedright\arraybackslash}p{4.0cm}
  >{\raggedright\arraybackslash}p{2.7cm}
  r r@{}}
\toprule
Campaign & Period & Analysers & Vertical design & Internal SPs & Rows\\
\midrule
1 & Jun.--Aug. 2024 & CRDS8, CRDS9, FTIR1, FTIR2 &
top, middle, bottom & 51 & 111\,210\\
2 & 16 Nov.--31 Dec. 2024 & CRDS8, CRDS9 &
top and bottom & 32 & 20\,632\\
\bottomrule
\end{tabular}
\endgroup
\end{table}

\subsection{Gas ratios}

The relative mixing ratios of \(\mathrm{CH_4}\) and \(\mathrm{NH_3}\) to
\(\mathrm{CO_2}\) were calculated for positive, valid concentrations:

\begin{align}
R_{\mathrm{CH_4:CO_2}} &=
100\,C_{\mathrm{CH_4}}/C_{\mathrm{CO_2}},\\
R_{\mathrm{NH_3:CO_2}} &=
100\,C_{\mathrm{NH_3}}/C_{\mathrm{CO_2}},\\
R_{\mathrm{NH_3:CH_4}} &=
100\,C_{\mathrm{NH_3}}/C_{\mathrm{CH_4}}.
\end{align}

The ratios are dimensionless and are reported as percentages. Where
inferential modelling of ratios is subsequently performed, their logarithmic
forms, \(\log(C_{\mathrm{CH_4}}/C_{\mathrm{CO_2}})\) and
\(\log(C_{\mathrm{NH_3}}/C_{\mathrm{CO_2}})\), are preferred to avoid the
asymmetry of untransformed ratios.

\subsection{Evaluation of sampling-density reduction}

Daily means at each location were first calculated to provide independent
blocks for descriptive means and 95\% confidence intervals. For direct
comparison of sampling configurations, measurements were averaged into
two-hour blocks. A block was retained only when all 51 internal locations
were represented, yielding 743 complete blocks in the recovered dataset.

Seven vertical configurations were reconstructed from the Campaign 1 data:
top only, middle only, bottom only, top plus middle, middle plus bottom, top
plus bottom, and the full three-height design. The full 51-location mean was
used as the reference estimate within each block. For a reduced configuration
\(k\), the relative deviation was calculated as

\begin{equation}
RE_{k,t}=100\,
\frac{\overline{C}_{k,t}-\overline{C}_{51,t}}
     {\overline{C}_{51,t}}.
\end{equation}

The top-and-bottom configuration was selected a priori as the operational
reduced design because it retained both ends of the observed vertical
gradient. Its agreement with the full design was quantified using mean bias,
relative bias, RMSE, coefficient of determination, and Lin's CCC. Spatial
heterogeneity within each block was expressed as the standard deviation of
log-transformed location concentrations.

Following \citet{Chen2026}, the temporal informativeness of every location was
described using Shannon entropy. For each gas and ratio, location means from
complete two-hour blocks were divided into four equal-width bins. The
empirical probabilities \(p_b\) were used to calculate
\begin{equation}
H^*=-\frac{\sum_{b=1}^{4}p_b\log_2(p_b)}{\log_2(4)}.
\end{equation}
This was an empirical adaptation of the cross-case CFD entropy method: the
present cases were observed two-hour blocks rather than simulated ventilation
scenarios. Sensitivity to temporal aggregation was assessed at 2, 4, 8, and
24 h resolutions.

\subsection{Multiscale persistence and reference-location analysis}

To determine whether location differences diminished through time, location
means were calculated at 1, 2, 4, 8, 12, 24, 72, and 168 h resolutions. Within
each contemporaneous block, every location was compared with the spatial mean
of the measured network. The decay in median absolute relative error was
described by a power-law regression of log error against log averaging
duration. Daily location-rank stability was quantified using pairwise Spearman
correlations and Kendall's coefficient of concordance, \(W\).

In addition, descriptive statistics were calculated explicitly at hourly,
daily, weekly, and whole-campaign scales for all three concentrations and the
three ratios. For each location and scale these comprised observation count,
mean, median, minimum, 5th and 95th percentiles, maximum, standard deviation
(SD), median absolute deviation (MAD), and coefficient of variation (CV).
Whole-campaign statistics were calculated directly from all valid observations
rather than by averaging weekly summaries. Repeated-location Friedman tests
were used to assess whether SD, MAD, and CV differed among the four reporting
scales.

Variance was partitioned using linear mixed-effects models of log-transformed
concentrations and ratios, with random intercepts for location and date and
fixed effects for campaign, height, and hour. Campaign-specific generalised
additive models included cyclic hour-of-day smooths, longer-term date smooths,
and location random effects. PCA and correlation-distance hierarchical
clustering were applied to hourly location relative-error profiles. Cluster
number was selected by the greatest mean silhouette width.

Reference-location selection was defined before statistical ranking. Location
40 was excluded because of confirmed sampling-line leakage and location 19
because of local fan influence. A conservative physical sensitivity screen
also excluded the fan corridors containing locations 19--21 and 43--45. Only
locations measured in both campaigns were eligible. Within each gas and
campaign, candidates were ranked for absolute bias, MAE, RMSE, temporal
stability, Spearman association and four-bin mutual information with the
contemporaneous spatial mean. Equally weighted percentile scores were then
averaged across \(\mathrm{CO_2}\), \(\mathrm{CH_4}\), and
\(\mathrm{NH_3}\); ratios were retained as a sensitivity analysis.

Predictive pattern recognition compared a linear model, cyclic GAM, and
regression tree for hourly location relative error. Five-fold cross-validation
was grouped by complete date, preventing observations from the same date from
appearing in both training and test data. Change-point screening of daily barn
means and spatial CV was exploratory and was not interpreted causally without
corresponding operational records.

All cleaning, calculations, and figures were produced in R version 4.5.3
(R Foundation for Statistical Computing, Vienna, Austria). Data manipulation
was conducted using \texttt{data.table}; mixed-effects and additive models
were fitted using \texttt{lme4} and \texttt{mgcv}, respectively; clustering
and silhouette analysis used \texttt{cluster}; regression trees used
\texttt{rpart}; and figures were produced using \texttt{ggplot2}. The complete
analysis code and intermediate summary tables were retained to permit the
sampling configurations and ranking criteria to be reproduced. Top, middle,
and bottom locations were displayed in orange, green, and light blue,
respectively.

\section{Results}

\subsection{High-resolution concentration patterns}

Marked spatial and vertical variation was detected during Campaign 1
(Fig.~\ref{fig:locations}). Across internal locations, the mean corrected
\(\mathrm{CO_2}\) concentrations were approximately \SI{616}{ppm} at the top,
\SI{611}{ppm} at the middle, and \SI{609}{ppm} at the bottom.
Corresponding \(\mathrm{CH_4}\) means were 18.84, 18.51, and
\SI{18.38}{ppm}, while \(\mathrm{NH_3}\) means were 2.05, 2.04, and
\SI{2.02}{ppm}. The direction and magnitude of the vertical gradient were not
uniform among horizontal positions. Particularly distinct values occurred at
locations 19 and 40, supporting their exclusion from the subsequent
operational design.

\begin{figure}[htbp]
\centering
\includegraphics[width=\textwidth]{Fig2_june_august_spatial_means_with_CV_v03.png}
\caption{Campaign 1 corrected concentration and gas-ratio means across the 17
horizontal positions. Symbol size represents the location-specific temporal
coefficient of variation, capped at 100\% for display. Orange, green, and
light-blue symbols denote top, middle, and bottom locations, respectively.
The south-background location \(s\) was excluded.}
\label{fig:locations}
\end{figure}

\subsection{Effect of vertical sampling configuration}

Use of one height alone produced systematic deviations from the full-density
mean (Fig.~\ref{fig:configurations}; Table~\ref{tab:configurations}).
Top-only sampling overestimated \(\mathrm{CO_2}\), \(\mathrm{CH_4}\), and
\(\mathrm{NH_3}\) by mean values of 0.68\%, 1.01\%, and 1.18\%,
respectively. Bottom-only sampling underestimated the same gases by 0.50\%,
0.78\%, and 1.40\%. The middle-only configuration was closer to the full mean,
but removed the vertical extremes required for evaluating heterogeneity.

Combining top and bottom locations substantially reduced configuration bias.
The mean absolute relative errors were 0.33\% for \(\mathrm{CO_2}\), 0.95\%
for \(\mathrm{CH_4}\), and 0.77\% for \(\mathrm{NH_3}\). Their 95th
percentiles were 0.85, 2.52, and 1.92\%, respectively.

\begin{figure}[htbp]
\centering
\includegraphics[width=\textwidth]{Fig3_campaign1_sampling_configuration_deviation.png}
\caption{Deviation of six reduced vertical sampling configurations from the
full 51-location mean in 743 complete two-hour blocks. Boxes show the
interquartile range and centre lines show medians; whiskers extend to
1.5 times the interquartile range. Green and yellow bands indicate deviations
within \(\pm5\%\) and \(\pm10\%\), respectively. Configurations labelled top,
middle, or bottom contain 17 locations; paired-height configurations contain
34 locations. Outliers are omitted visually but retained in all summaries.}
\label{fig:configurations}
\end{figure}

\begin{table}[htbp]
\centering
\caption{Performance of selected reduced configurations relative to the full
51-location mean. Values are calculated across 743 complete two-hour blocks.
MAD denotes median absolute deviation and \(P_{95}\) denotes the 95th
percentile of absolute deviation.}
\label{tab:configurations}
\begin{tabular}{llrrr}
\toprule
Gas & Configuration & Mean RE (\%) & MAD (\%) & \(P_{95}\) (\%)\\
\midrule
\(\mathrm{CO_2}\) & Top only & 0.68 & 0.82 & 3.71\\
 & Middle only & -0.18 & 0.52 & 1.70\\
 & Bottom only & -0.50 & 0.78 & 3.28\\
 & Top + bottom & 0.09 & 0.26 & 0.85\\
\(\mathrm{CH_4}\) & Top only & 1.01 & 2.42 & 9.75\\
 & Middle only & -0.23 & 1.46 & 5.03\\
 & Bottom only & -0.78 & 2.17 & 8.37\\
 & Top + bottom & 0.12 & 0.73 & 2.52\\
\(\mathrm{NH_3}\) & Top only & 1.18 & 2.34 & 9.67\\
 & Middle only & 0.22 & 1.28 & 3.84\\
 & Bottom only & -1.40 & 2.02 & 8.04\\
 & Top + bottom & -0.11 & 0.64 & 1.92\\
\bottomrule
\end{tabular}
\end{table}

\subsection{Entropy and temporal-resolution sensitivity}

The most temporally informative location depended on the response
(Fig.~\ref{fig:entropy}). The highest normalised entropy occurred at locations
4, 7, and 15 for \(\mathrm{CO_2}\), \(\mathrm{CH_4}\), and
\(\mathrm{NH_3}\), respectively; no single location was universally
informative. Increasing the averaging interval from 2 to 24 h reduced the
median spatial CV from 7.73 to 4.60\% for \(\mathrm{CO_2}\), from 19.11 to
9.56\% for \(\mathrm{CH_4}\), and from 22.88 to 20.11\% for
\(\mathrm{NH_3}\). The persistent ammonia CV indicated that its
heterogeneity was not primarily a short-timescale artefact.

\begin{figure}[htbp]
\centering
\includegraphics[width=\textwidth]{Fig6_june_august_entropy_v03.png}
\caption{Normalised four-bin Shannon entropy of corrected gas concentrations
and gas ratios at each Campaign 1 location. Entropy was calculated across 743
complete two-hour blocks and normalised by \(\log_2(4)\). Numbers identify
sampling locations. Higher values indicate broader and more even occupation
of the four empirical bins, not necessarily a greater mean concentration.}
\label{fig:entropy}
\end{figure}

\FloatBarrier
\subsection{Concentrations and variation over time}

Barn-wide daily means varied substantially within both analysed periods. Across
the 89 represented Campaign 1 dates, the daily mean (mean across dates) was
\SI{607.6}{ppm} for \(\mathrm{CO_2}\), \SI{18.14}{ppm} for
\(\mathrm{CH_4}\), and \SI{2.08}{ppm} for \(\mathrm{NH_3}\). Their daily
means ranged from 506.0 to \SI{684.3}{ppm}, from 9.87 to
\SI{22.92}{ppm}, and from 1.25 to \SI{3.12}{ppm}, respectively. The
between-date CV was 6.36\% for \(\mathrm{CO_2}\), 14.31\% for
\(\mathrm{CH_4}\), and 16.89\% for \(\mathrm{NH_3}\). Thus, ammonia
showed the largest relative day-to-day variation during the dense summer
period.

Across the 47 represented Campaign 2 dates, the corresponding daily means were
\SI{636.9}{ppm}, \SI{18.40}{ppm}, and \SI{0.91}{ppm}. Daily means ranged
from 490.1 to \SI{802.7}{ppm} for \(\mathrm{CO_2}\), from 7.33 to
\SI{31.24}{ppm} for \(\mathrm{CH_4}\), and from 0.53 to
\SI{1.64}{ppm} for \(\mathrm{NH_3}\). Between-date CVs increased to
13.31\%, 36.16\%, and 27.83\%, respectively. Carbon dioxide and methane
reached their lowest and highest daily barn means on the same dates within
each campaign, whereas the ammonia extrema occurred on different dates. This
temporal decoupling was consistent with the gas-specific spatial and ratio
behaviour. These summaries describe the June--August component of Campaign 1
and Campaign 2; the recovered October component will require the same
calculation before an autumn comparison is reported.

\subsection{Persistence, convergence, and pattern recognition}

Spatial differences became smaller but did not disappear with longer
averaging (Fig.~\ref{fig:convergence}). In Campaign 1, median absolute location
error decreased from 6.43 to 3.12\% for \(\mathrm{CO_2}\), from 15.65 to
4.86\% for \(\mathrm{CH_4}\), and from 17.15 to 12.01\% for
\(\mathrm{NH_3}\) between hourly and weekly aggregation. Campaign 2 errors
decreased from 5.73 to 3.65\%, 17.91 to 10.24\%, and 17.34 to 15.37\%,
respectively. All power-law exponents were negative, although convergence was
weakest for \(\mathrm{NH_3}\).

The explicit four-scale summaries supported this distinction
(Fig.~\ref{fig:four-scale-cv}). Median within-location CV decreased from hourly
to weekly resolution in both campaigns: from 16.33 to 4.40\% and from 16.38 to
6.39\% for \(\mathrm{CO_2}\); from 38.37 to 9.45\% and from 48.17 to
19.79\% for \(\mathrm{CH_4}\); and from 31.47 to 18.32\% and from 28.60 to
13.25\% for \(\mathrm{NH_3}\) in Campaigns 1 and 2, respectively. The same
general attenuation occurred for all three ratios. Friedman tests showed scale
effects on SD, MAD, and CV for every concentration and ratio in both campaigns
(all \(P<10^{-12}\)). In contrast, whole-campaign CV calculated from all raw
observations remained high because it retained the complete seasonal and
diurnal range; it is therefore a descriptor of total campaign variability, not
an additional temporal-averaging step.

\begin{figure}[htbp]
\centering
\includegraphics[width=0.94\textwidth]{Fig_01_temporal_convergence.png}
\caption{Convergence of median absolute location error as the averaging
duration increased from 1 to 168 h. Errors were calculated relative to the
contemporaneous measured-network mean. The logarithmic horizontal axis shows
that concentration fields became more representative with aggregation but
retained gas- and campaign-specific residual heterogeneity.}
\label{fig:convergence}
\end{figure}

\begin{figure}[htbp]
\centering
\includegraphics[width=0.94\textwidth]{Fig_04b_hour_day_week_campaign_CV.png}
\caption{Median within-location coefficient of variation (CV) for
\(\mathrm{CO_2}\), \(\mathrm{CH_4}\), \(\mathrm{NH_3}\), and their three
ratios at hourly, daily, weekly, and whole-campaign reporting scales. Hourly,
daily, and weekly values quantify variability among block means, whereas the
whole-campaign value was calculated directly from every valid observation and
therefore preserves the full seasonal and diurnal range. Consequently, the
whole-campaign bar should not be interpreted as a continuation of the temporal
averaging sequence.}
\label{fig:four-scale-cv}
\end{figure}

Daily rank stability was moderate in Campaign 1: Kendall's \(W\) was 0.433 for
\(\mathrm{CO_2}\), 0.447 for \(\mathrm{CH_4}\), and 0.523 for
\(\mathrm{NH_3}\). Location accounted for 4.4\% of log-\(\mathrm{CO_2}\)
variance, 3.7\% of log-\(\mathrm{CH_4}\) variance, and 21.2\% of
log-\(\mathrm{NH_3}\) variance. The corresponding location fractions were
4.9\% for \(\mathrm{CH_4}/\mathrm{CO_2}\), 22.9\% for
\(\mathrm{NH_3}/\mathrm{CO_2}\), and 20.6\% for
\(\mathrm{NH_3}/\mathrm{CH_4}\). PCA and clustering did not identify one
universal grouping for all gases, indicating response-specific spatial
patterns.

Regression trees outperformed linear and GAM benchmarks in date-grouped
cross-validation, but predictive performance remained moderate. Mean test
\(R^2\) was 0.205 for \(\mathrm{CO_2}\), 0.121 for \(\mathrm{CH_4}\), and
0.409 for \(\mathrm{NH_3}\). Location and clock time therefore captured only
part of hourly spatial error, leaving substantial variation for ventilation,
weather, animal, and operational covariates.

\subsection{Pragmatic reference-location selection}

Location 25 ranked first among physically eligible cross-campaign candidates
(Fig.~\ref{fig:reference}). It was outside the conservative fan corridors, was
measured in both campaigns, and remained first when ratio behaviour was added
as a sensitivity criterion. Its hourly biases in Campaigns 1 and 2 were
-1.07 and +0.62\% for \(\mathrm{CO_2}\), +3.00 and +0.88\% for
\(\mathrm{CH_4}\), and -2.09 and +7.25\% for \(\mathrm{NH_3}\),
respectively. It was therefore the best single-point compromise rather than a
universally unbiased substitute for spatial averaging.

\begin{figure}[htbp]
\centering
\includegraphics[width=0.80\textwidth]{Fig_03_reference_location_ranking.png}
\caption{Leading cross-campaign reference-location candidates. The primary
score equally combines concentration bias, absolute error, temporal stability,
correlation, and mutual information after the physical exclusion criteria.
The worst-gas score penalises candidates whose performance was poor for any
one gas. Fan-corridor candidates were retained only in the sensitivity ranking.}
\label{fig:reference}
\end{figure}

\begin{figure}[htbp]
\centering
\includegraphics[width=\textwidth]{Fig_05_reference_location_floorplan.png}
\caption{Physical screening of the statistically selected reference candidate.
Location 25 is marked in blue. Red crosses identify the conservatively excluded
fan corridors containing locations 19--21 and 43--45; the black cross marks
location 40, which was excluded because of confirmed sampling-line leakage.
The exact fan axes should be confirmed against the operational installation
before final engineering acceptance.}
\label{fig:reference-map}
\end{figure}

\subsection{Agreement between full and reduced designs}

The top-and-bottom mean closely reproduced the full 51-location estimate
throughout Campaign 1 (Fig.~\ref{fig:agreement}). Lin's CCC was 0.999 for all
three gases. Mean relative biases were 0.09\% for \(\mathrm{CO_2}\), 0.12\%
for \(\mathrm{CH_4}\), and -0.11\% for \(\mathrm{NH_3}\). The corresponding
RMSE values were \SI{2.69}{ppm}, \SI{0.23}{ppm}, and \SI{0.019}{ppm}.
Agreement for the spatial heterogeneity metric was lower than agreement for
the mean but remained strong, with CCC values between 0.953 and 0.968.

\begin{figure}[htbp]
\centering
\includegraphics[width=\textwidth]{Fig5_june_august_agreement_v03.png}
\caption{Agreement between the full 51-location mean and the reduced
34-location top-and-bottom mean in 743 complete two-hour blocks. Solid lines
represent exact 1:1 agreement. CCC denotes Lin's concordance correlation
coefficient, and bias is the mean relative difference from the full-density
estimate.}
\label{fig:agreement}
\end{figure}

\subsection{Longer-term application of the reduced design}

The reduced-height principle was applied from mid-November to the end of
December 2024 in Campaign 2 (Fig.~\ref{fig:campaign2}). Mean concentrations
remained higher at top than at bottom locations: \(\mathrm{CO_2}\) averaged
approximately 640 and \SI{625}{ppm}, \(\mathrm{CH_4}\) averaged 18.6 and
\SI{17.4}{ppm}, and \(\mathrm{NH_3}\) averaged 0.96 and \SI{0.83}{ppm} at
the top and bottom, respectively. The mean \(\mathrm{CH_4}/\mathrm{CO_2}\)
ratios were 2.75\% and 2.64\%, and the mean
\(\mathrm{NH_3}/\mathrm{CO_2}\) ratios were 0.149\% and 0.131\% at the top
and bottom.

\begin{figure}[htbp]
\centering
\includegraphics[width=\textwidth]{Fig5_campaign2_concentrations_reduced_design.png}
\caption{Campaign 2 concentrations measured with the reduced top-and-bottom
network. Points denote means of daily location means and vertical bars denote
95\% confidence intervals calculated across daily blocks. Orange and
light-blue symbols denote top and bottom locations, respectively.}
\label{fig:campaign2}
\end{figure}

\section{Discussion}

\subsection{Gas-specific spatial structure}

The measurements did not support homogeneous concentration fields within the
investigated barn. Vertical differences were systematic at the network level,
but their magnitude varied among gases and horizontal positions. This
gas-specific structure is consistent with earlier field observations showing
height-dependent \(\mathrm{NH_3}\), \(\mathrm{CO_2}\), and tracer-gas
mixing ratios \citep{Mendes2015}, and with the incomplete agreement among
three measurement heights reported for an open dairy barn
\citep{Durso2024}. The present study extends that evidence by resolving the
vertical comparison at 17 horizontal positions rather than at one or a few
profiles.

Horizontal coverage remained important even when the vertical design was
reduced. This agrees with \citet{Durso2022}, for whom both the number and
spacing of horizontal locations affected the estimated concentration
distribution. The current results further show why vertical and horizontal
design cannot be considered independently: omission of one height may be
acceptable for the spatial mean only when the retained heights remain
distributed across the barn. The variance partitioning also demonstrated that
this requirement was response-specific. Location explained a substantially
larger fraction of variation in \(\mathrm{NH_3}\) and its ratios than in
\(\mathrm{CO_2}\) or \(\mathrm{CH_4}\), indicating that a network adequate
for one gas cannot automatically be assumed adequate for another.

Several mechanisms are compatible with the observed patterns, although they
were not measured directly in these campaigns. Carbon dioxide was dominated
by a spatially distributed animal source, whereas methane included localised
enteric release and ammonia was linked to manure-contaminated surfaces and
gas--surface exchange. Wind direction, opening configuration, buoyancy, and
fan operation could then transport or dilute these sources differently.
External wind conditions have previously affected internal concentrations and
air exchange in a naturally ventilated dairy building \citep{Saha2013}, but
the absence of concurrent airflow-vector measurements prevents individual
features in the present concentration maps from being attributed causally to
a specific opening, fan, or flow path.

\subsection{Information contained in the gas ratios}

The ratios were not interpreted as source signatures in isolation. Instead,
they were used as diagnostics of whether gases varied proportionally across
space and time. The comparatively small location fraction for
\(\mathrm{CH_4}/\mathrm{CO_2}\) suggests that part of the spatial variation
in these two animal-associated gases was shared. In contrast, the larger
location fractions for \(\mathrm{NH_3}/\mathrm{CO_2}\) and
\(\mathrm{NH_3}/\mathrm{CH_4}\) indicate that ammonia retained spatial
behaviour not removed by normalisation to either carbon gas. This finding is
consistent with the spatial variability of ammonia-to-tracer mixing ratios
reported by \citet{Mendes2015} and reinforces the need to evaluate ratio
stability rather than assume co-dispersion.

Ratio interpretation remains conditional on the source distributions and
background treatment. No background-corrected enhancement regression was
applied, and the south-background location was excluded from the internal
heterogeneity analysis. Consequently, the ratios quantify proportionality of
the measured internal concentrations; they do not by themselves distinguish
animal, manure, or ventilation contributions and should not be interpreted as
emission factors.

\subsection{Temporal aggregation does not replace spatial coverage}

\subsubsection{Temporal concentration dynamics}

The daily barn means show that temporal variability was not merely scatter
among sampling locations. Concentrations changed coherently at the barn scale,
but the relative magnitude and timing differed among gases. The coincident
daily extrema of \(\mathrm{CO_2}\) and \(\mathrm{CH_4}\) are compatible
with shared effects of animal occupancy, activity, and ventilation, while the
different timing and larger relative variation of \(\mathrm{NH_3}\) are
compatible with additional control by manure-surface conditions and
gas--surface exchange. These are supported interpretations rather than causal
tests because animal activity, surface state, airflow, and meteorology have not
yet been included together in the fitted models.

Campaign 2 had wider daily relative variation for all three gases, but its
absolute \(\mathrm{NH_3}\) level was lower than in the analysed summer part
of Campaign 1. This difference cannot be assigned to season alone: campaign,
network density, omitted locations, analyser composition, ventilation regime,
and unmodelled barn operation changed together. The recovered October dense
period offers the more defensible first test of seasonal persistence because
it retains the Campaign 1 spatial design. Its integration with on-site mast
and regional meteorological data should determine whether changes in temporal
concentration level coincide with changes in spatial ranks or only with
barn-wide dilution.

Aggregation reduced median absolute location error for all three gases, but
the rate and remaining error differed. The stronger persistence of
\(\mathrm{NH_3}\) relative to \(\mathrm{CO_2}\) and
\(\mathrm{CH_4}\) was consistent with its larger location variance fraction.
Thus, short-lived fluctuations were averaged out while stable or repeatedly
occurring spatial contrasts remained. The result clarifies two different
meanings of representativeness: a longer average can improve the precision of
a location mean without removing its systematic displacement from the
contemporaneous network mean.

Whole-campaign variability should also not be read as the final point of a
monotonic averaging sequence. The campaign CV was calculated from all valid
observations and therefore retained diurnal and seasonal variation, whereas
the daily and weekly CVs described variation among aggregated block means.
The high whole-campaign CV therefore reflects a broader temporal domain, not a
failure of weekly averaging. Sequential sampling adds a further limitation:
rapid changes during a switching cycle can be aliased as spatial differences,
and adsorption or delayed response is particularly relevant for ammonia
\citep{Rom2010}. Complete two-hour blocks reduced, but could not eliminate,
this temporal--spatial ambiguity.

\subsection{Reduction of the vertical network}

The middle-only design reproduced the dense-network mean comparatively well,
but it suppressed the vertical extremes that defined the heterogeneity under
investigation. The top-and-bottom design retained these extremes and allowed
their opposing deviations to balance in the spatial mean. Its near-unity CCC
and small bias demonstrate agreement with the 51-location reference estimate
for the observed two-hour blocks; they do not prove that every local plume or
measure of spatial spread was preserved.

The finding is consistent with the broader observation that sensor placement
is not interchangeable with sensor number. Different arrangements at the same
sampling density produced different airflow estimates in a naturally
ventilated mock-up building \citep{DeVogeleer2017}. Similarly, vertically
composite concentration and velocity measurements were more representative
than a single-height approach in a wind-tunnel study
\citep{Janke2020}. Those studies concern airflow or emission estimation,
whereas the present validation concerns internal concentration means. The
agreement across these different outcomes supports vertically distributed
sampling as a design principle, but does not make their numerical errors
directly comparable.

Campaign 2 demonstrated that a top-and-bottom network could be operated over
a longer late-autumn and winter period. It did not provide an independent
51-versus-34 validation because no simultaneous dense reference was present
and two site-specific locations had also been removed. The evidence therefore
supports a reduction from 51 to 34 locations within Campaign 1 and an
operational 32-location application in Campaign 2, rather than a universal
32-location optimum.

\subsection{Information-based ranking and the single-point question}

The entropy analysis adapted the information-based placement principle
described by \citet{Chen2026}. In that study,
sensor informativeness was evaluated across CFD scenarios using local mean age
of air and related flow variables. Here, entropy described temporal diversity
of observed gas concentrations at individual locations. It was therefore a
diagnostic component rather than a direct implementation of the CFD-informed
optimisation algorithm, and the reduced vertical design had been defined from
the configuration comparison rather than selected by entropy alone.

Likewise, the grouped-date predictive models showed only moderate out-of-date
performance. This is scientifically informative: location, height, and clock
time captured repeatable structure, but much of the hourly relative error
remained unexplained. Weather, fan status, animal distribution, and barn
microclimate are plausible omitted predictors. The models should therefore be
viewed as pattern-recognition benchmarks, not operational concentration
predictors.

The recovered autumn dense-network period is particularly important for the
next model revision because it retains the same 51-location geometry and
four-analyser deployment used in the summer period. It can therefore test
whether location ranks, entropy, and reduced-design agreement persist under a
different meteorological regime without simultaneously changing the spatial
network. Wind speed should be allowed a non-linear response, and direction
should be represented by sine and cosine components or pre-specified wind
sectors rather than as a linear 0--360\(^\circ\) covariate. Temperature and
relative humidity are physically justified covariates, while precipitation is
more appropriately evaluated as an event indicator and lagged modifier than
assumed to exert an immediate independent effect. These analyses have not yet
been used to revise the numerical results reported above.

Location 25 was the best compromise among candidates that passed the
predefined physical screen and were available in both campaigns. Its small
carbon-gas biases support its use as a pragmatic single-point approximation
within this barn. However, the \(\mathrm{NH_3}\) bias changed to +7.25\% in
Campaign 2, demonstrating residual response- and campaign-dependence. A point
that approximates the internal concentration mean is also not automatically
suitable for emission calculation because flux depends jointly on
concentration and airflow across a defined control surface
\citep{Calvet2013,Ogink2013}. Location 25 should therefore be treated as a
testable site-specific candidate rather than a reference standard.

\subsection{Limitations and future research}

The principal limitations define the scope of inference. The summer and
autumn periods belong to Campaign 1 and retain the dense design, whereas the
late-autumn and winter period belongs to Campaign 2 and uses a reduced network
and different analyser deployment. Consequently, differences between
Campaigns 1 and 2 cannot be attributed causally to season alone. Comparisons
between the summer and autumn parts of Campaign 1 provide a less confounded
test of temporal persistence, but analyser and location remained partly
confounded because the four analysers were assigned to fixed spatial ranges.
The applied factors corrected mean FTIR--CRDS offsets but did not establish
traceability to an absolute reference or remove subsequent drift. Sequential
switching prevented an instantaneous three-dimensional snapshot, and the
51-location mean remained a dense-network reference estimate rather than a
volume-weighted true concentration.

Future validation should randomise or exchange analyser assignments. Sampling
line response should be quantified for the deployed tube lengths and flow
rates, and reduced and dense networks should be compared concurrently in
independent periods.
Meteorological, fan, animal-distribution, and barn-climate variables should be
aligned with the concentration blocks to test the proposed mechanisms. Wind
measured at the on-site \SI{10}{m} mast should be treated as the primary
external-flow observation where available. Hourly DWD station observations of
wind, air temperature, relative humidity, and precipitation can provide a
quality-controlled regional context and support gap and sensitivity analyses,
but should not be treated as measurements of the air approaching the barn
after modification by nearby buildings, vegetation, and terrain.
CFD could then be evaluated against the empirical spatial patterns and used to
examine unmeasured airflow paths, but such modelling would constitute a new
validation study rather than evidence produced here. Finally, emission
applications require simultaneous assessment of representative airflow and
concentration at an appropriate control surface; the present concentration
ranking alone is insufficient for that purpose.

\section{Conclusions}

Concentrations in the investigated naturally ventilated dairy barn remained
heterogeneous across both height and horizontal position, and the persistence
of that structure was gas-specific. Temporal aggregation reduced short-term
location error but did not remove the comparatively strong spatial component
of ammonia and its ratios. A measurement period should therefore not be
lengthened as a substitute for spatial coverage without first demonstrating
that the remaining location bias is acceptable for the intended gas and
outcome.

The dense campaign showed that a distributed top-and-bottom configuration
could preserve the internal \(\mathrm{CO_2}\), \(\mathrm{CH_4}\), and
\(\mathrm{NH_3}\) reference means while retaining vertical contrast and
reducing sampling effort. Dense initial mapping followed by explicit,
response-specific configuration testing is therefore supported as a practical
design strategy. A single physically screened location could provide a useful
site-specific approximation, but it could not replace spatial averaging or be
assumed representative of emission flux. Independent seasonal validation with
concurrent airflow and operational measurements remains necessary before the
reduced configuration or single-point candidate is transferred beyond this
barn.

\section*{Acknowledgements}

The technical and animal-care staff of the experimental dairy facility are
gratefully acknowledged. \textit{[Funding project names, grant numbers, and
specific technical contributors are to be confirmed before submission.]}

\section*{Author contributions}

Harsh Sahu: conceptualisation, methodology, investigation, data curation,
formal analysis, visualisation, and writing--original draft. Long Chen:
methodology, formal analysis, and writing--review and editing. Thomas Amon:
conceptualisation, supervision, funding acquisition, and writing--review and
editing. Tariq Mehmood: investigation and writing--review and editing. Aditya
Rawat: investigation and writing--review and editing. David Janke:
conceptualisation, methodology, supervision, project administration, and
writing--review and editing.

\section*{Declaration of competing interest}

The authors declare that they have no known competing financial interests or
personal relationships that could have appeared to influence the work
reported in this paper.

\section*{Data availability}

The processed data and analysis code will be made available in a public
repository upon acceptance. \textit{[Repository and DOI to be added.]}

\bibliographystyle{apalike}
\bibliography{references}

\end{document}
